\section{Conclusion and Future Work}
\label{sec:conclusion}

In this work, we presented a rigorous deconstruction of the ``Layer-Averaging Collapse'' (rank-1 collapse) claim in dynamic model merging. Guided by a critical methodological perspective, we exposed that prior assertions of theoretical collapse are artifacts of over-simplified, linear representation-space sandboxes. 

By conducting a rigorous physical empirical evaluation on Split-MNIST digit subsets using deep neural network backbones (DeepMLP-12 and TinyCNN-4) across task suites of varying semantic conflict, we showed that learned dynamic routing trajectories occupy a multi-dimensional subspace. Specifically, the SVD Collinearity Ratio drops to $0.4987 \pm 0.08$ (on DeepMLP-12) and $0.5673 \pm 0.03$ (on TinyCNN-4) under Cross-Domain physical task conflict. Probing physical routing matrices with inter-layer cosine similarity heatmaps revealed the emergence of structured, depth-specialized routing blocks aligned with the network's hierarchical feature extraction. Furthermore, our proposed Bounded Sigmoid (BSigmoid) dynamic routing pipeline provides an expressive weight-space merging framework, while highlighting the vital role of optimizer regularization and parameter-variance trade-offs under few-shot splits.

Our study serves as a critical call to action: we must move away from toy linear representation-space sandboxes and establish more rigorous physical empirical evaluation protocols to properly understand and leverage weight-space dynamic model merging. Future work will scale physical layer-wise routing and the SVD collinearity audit to larger, over-parameterized architectures (such as Vision Transformers and Large Language Models) and more semantically diverse datasets (such as CIFAR-10, SVHN, or ImageNet). When scaling from highly localized 2D handwritten digits to high-dimensional natural images with complex textures, rotations, and lighting scales, we hypothesize that the geometric routing dynamics will exhibit even richer multi-dimensional subspace trajectories. As hierarchical representations in larger backbones specialize to capture complex natural semantics at different depths (e.g., low-level edge detectors in early layers versus task-specific object classes in deeper layers), the SVD collinearity ratio is expected to drop further, showing even more pronounced block-diagonal specialization patterns. Verifying these dynamics on complex natural image domains remains a prioritized path for future empirical validation.
